\section*{Capítulo \ref{ch:g9:natural-risks} --- Conviver com os riscos naturais}

\begin{solution}{exo:g9:natural-risks:1}
(a) Ameaça. (b) \omterm{def:g9:natural-risks:vulnerability}{Vulnerabilidade}. (c) \omterm{def:g9:natural-risks:vulnerability}{Vulnerabilidade}. (d) Ameaça. (e) As
duas: os fluxos de lama são a ameaça; construir a cidade no caminho
deles a torna vulnerável.
\end{solution}

\begin{solution}{exo:g9:natural-risks:2}
(a) Falso: onde ninguém mora, nada pode ser atingido, então o \omterm{def:g9:natural-risks:risk}{risco} é
baixo. (b) Verdadeiro. (c) Verdadeiro. (d) Falso: um alerta não muda o
\omterm{def:g4:volcanoes-earthquakes-first-look:volcano}{vulcão}; ele reduz a \omterm{def:g9:natural-risks:vulnerability}{vulnerabilidade} das pessoas, que podem sair a tempo.
\end{solution}

\begin{solution}{exo:g9:natural-risks:3}
(a) Sim, com dias de antecedência. (b) Sim, com horas ou dias de
antecedência. (c) Não: os \omterm{def:g4:volcanoes-earthquakes-first-look:earthquake}{terremotos} não podem ser previstos. (d) Sim:
depois que o \omterm{def:g7:earthquakes:tsunami}{tsunami} começou, ele leva horas para atravessar o \omterm{def:g2:water-on-earth:ocean}{oceano}, e
o alerta viaja mais depressa.
\end{solution}

\begin{solution}{exo:g9:natural-risks:4}
(a) O tremor do chão: enxames de pequenos \omterm{def:g4:volcanoes-earthquakes-first-look:earthquake}{terremotos} quando o \omterm{def:g7:volcanoes:magma}{magma} sobe.
(b) A inclinação da encosta: o \omterm{def:g4:volcanoes-earthquakes-first-look:volcano}{vulcão} que incha. (c) Como o chão se move,
em milímetros: o inchaço quando o \omterm{def:g7:volcanoes:magma}{magma} se acumula. (d) Os gases da
pluma: mais gás escapando quando o \omterm{def:g7:volcanoes:magma}{magma} sobe.
\end{solution}

\begin{solution}{exo:g9:natural-risks:5}
(a) $45 \div 2 = \qty{22.5}{km/h}$. (b) $4 \times 2 = \qty{8}{km}$, o
bastante para chegar a um terreno mais alto. (c) Ninguém viu a \omterm{def:g4:volcanoes-earthquakes-first-look:eruption}{erupção},
escondida pela \omterm{def:g5:natural-hazards:storm}{tempestade}, e ninguém os avisou: não havia meio rápido de
enviar a informação à cidade, nem plano de evacuação. (d) Por exemplo: um
plano de evacuação com caminhos até um terreno mais alto, treinado com
antecedência; uma corrente de alerta dos cientistas até a cidade; mais
instrumentos no \omterm{def:g4:volcanoes-earthquakes-first-look:volcano}{vulcão}.
\end{solution}

\begin{solution}{exo:g9:natural-risks:6}
(a) A \omterm{def:g5:water-shapes-the-land:flood}{inundação} \qty{10}{m} acima do \omterm{def:g1:landscapes-around-us:water}{rio}. (b) Cerca de \qty{0.47}{km}
para a \omterm{def:g5:water-shapes-the-land:flood}{inundação} de \qty{2}{m}, \qty{0.56}{km} para a de \qty{5}{m} e
\qty{0.67}{km} para a de \qty{10}{m} (leituras com diferença de até cerca
de \qty{0.05}{km} são aceitáveis). (c) Só a \omterm{def:g5:water-shapes-the-land:flood}{inundação} de \qty{10}{m}: o
chão está mais de 2 e de \qty{5}{m} acima do \omterm{def:g1:landscapes-around-us:water}{rio}. (d) Basta uma pequena
subida do \omterm{def:g1:landscapes-around-us:water}{rio} para alcançá-la, e as \omterm{def:g5:water-shapes-the-land:flood}{inundações} pequenas vêm com mais
frequência do que as grandes.
\end{solution}

\begin{solution}{exo:g9:natural-risks:7}
(a) Ameaça alta, \omterm{def:g9:natural-risks:vulnerability}{vulnerabilidade} baixa: \omterm{def:g9:natural-risks:risk}{risco} moderado. (b) Ameaça baixa,
\omterm{def:g9:natural-risks:vulnerability}{vulnerabilidade} alta: \omterm{def:g9:natural-risks:risk}{risco} moderado. (c) Ameaça alta, \omterm{def:g9:natural-risks:vulnerability}{vulnerabilidade}
alta: \omterm{def:g9:natural-risks:risk}{risco} alto. (d) Ameaça baixa, \omterm{def:g9:natural-risks:vulnerability}{vulnerabilidade} baixa: \omterm{def:g9:natural-risks:risk}{risco} baixo.
\end{solution}

\begin{solution}{exo:g9:natural-risks:8}
(a) C, depois B, depois A. (b) As travas formam triângulos, que não podem
se inclinar sem esticar ou esmagar uma trava: a estrutura mantém a sua
forma em vez de se inclinar. (c) Os isoladores cedem quando o chão se
move para a frente e para trás embaixo deles, de modo que muito menos do
movimento chega à escola.
\end{solution}

\begin{solution}{exo:g9:natural-risks:9}
(a) $1 \div 100 = 1~\%$ por ano. (b) A chance é a mesma todo ano,
aconteça o que tiver acontecido no ano anterior: a \omterm{def:g5:water-shapes-the-land:flood}{inundação} pode voltar
no ano que vem. (c) $300 \div 100 = 3$ \omterm{def:g5:water-shapes-the-land:flood}{inundações}, em média.
\end{solution}

\begin{solution}{exo:g9:natural-risks:10}
(a) $3600 \div 15 = 240$ leituras. (b) $240 \times 24 \times 365 =
\num{2102400} \approx 2.1 \times 10^6$ leituras. (c) $2400 \div 800 =
\qty{3}{h}$. (d) Nessas três horas, os centros de alerta podem avisar o
\omterm{def:g1:landscapes-around-us:water}{litoral}, e as pessoas podem ir para um terreno mais alto.
\end{solution}

\begin{solution}{exo:g9:natural-risks:11}
(a) $4 \times 3 \times 3.8 = \qty{45.6}{L}$. (b) $45.6 \div 1.5 = 30.4$,
então 31 garrafas. (c) Por exemplo: comida que se conserva, para comer
quando as lojas estiverem fechadas; uma lanterna, se a eletricidade for
cortada; um rádio a pilha ou a manivela, para ouvir os alertas; um kit de
primeiros socorros, para tratar ferimentos pequenos; cobertores quentes.
(d) Para que ele possa ser levado na hora, se a família tiver de sair
depressa.
\end{solution}

\begin{solution}{exo:g9:natural-risks:12}
Por exemplo: (1) desenhar um \omterm{def:g9:natural-risks:hazard-map}{mapa de ameaças} dos fluxos de lama, \omterm{def:g1:landscapes-around-us:landform}{vale} por
\omterm{def:g1:landscapes-around-us:landform}{vale} (lado da ameaça); (2) instalar \omterm{def:g7:earthquakes:seismometer}{sismômetros}, inclinômetros,
receptores GPS e medidores de gases no \omterm{def:g4:volcanoes-earthquakes-first-look:volcano}{vulcão} (lado da ameaça); (3)
montar uma corrente rápida de alerta do observatório até a cidade, com
sirenes (lado da \omterm{def:g9:natural-risks:vulnerability}{vulnerabilidade}); (4) planejar e sinalizar caminhos até
um terreno mais alto, e treiná-los com simulados (lado da
\omterm{def:g9:natural-risks:vulnerability}{vulnerabilidade}); (5) não construir novas casas nem escolas no caminho
dos fluxos de lama (lado da \omterm{def:g9:natural-risks:vulnerability}{vulnerabilidade}).
\end{solution}

\begin{solution}{pb:g9:natural-risks:1}
\textbf{1.} O \omterm{def:g1:landscapes-around-us:water}{rio} está a \qty{143}{m} ao lado de C, a \qty{137}{m} ao
lado de A e a \qty{126}{m} ao lado de B; no \omterm{def:g2:drawing-to-map:map}{mapa}, C fica a oeste e B a
leste, então o \omterm{def:g1:landscapes-around-us:water}{rio} corre para leste.

\textbf{2.} A: $137 - 137 = \qty{0}{m}$; B: $276 - 126 = \qty{150}{m}$;
C: $236 - 143 = \qty{93}{m}$.

\textbf{3.} Só A, que fica a menos de \qty{5}{m} acima do \omterm{def:g1:landscapes-around-us:water}{rio}.

\textbf{4.} Cerca de \qty{0.57}{km} (entre 0.5 e \qty{0.65}{km}).

\textbf{5.} A subida em \qty{100}{m} é metade da subida em
\qty{200}{m}: A $1.4 \div 2 = 0.7~\%$; B $83 \div 2 = 41.5~\%$;
C $20 \div 2 = 10~\%$.

\textbf{6.} B: a sua encosta é muito íngreme, e \omterm{def:g3:stones-pebbles-sand:rock}{rochas} e terra soltas
podem deslizar por uma encosta íngreme, sobretudo depois de uma \omterm{def:g4:clouds-rain-snow:precipitation}{chuva}
forte.

\textbf{7.} A: cerca de \qty{1}{km}; B: quase sobre a \omterm{def:g7:earthquakes:fault}{falha} (a menos de
\qty{0.05}{km}); C: cerca de \qty{1.3}{km}.

\textbf{8.} Se a \omterm{def:g7:earthquakes:fault}{falha} se move, o chão se rompe e se desloca ao longo
dela, o que pode rasgar um prédio ao meio; e um \omterm{def:g4:volcanoes-earthquakes-first-look:earthquake}{terremoto} na \omterm{def:g7:earthquakes:fault}{falha}
sacudiria o chão ali com muita força.

\textbf{9.}

\begin{center}
{\small
\begin{tabular}{|c|c|c|c|}
\hline
terreno & na zona de \omterm{def:g5:water-shapes-the-land:flood}{inundação}? & encosta íngreme? & sobre a \omterm{def:g7:earthquakes:fault}{falha}? \\
\hline
A & sim & não & não \\
B & não & sim & sim \\
C & não & não & não \\
\hline
\end{tabular}}
\end{center}

\textbf{10.} O terreno C: é o único terreno com três ``não'':
\qty{93}{m} acima do \omterm{def:g1:landscapes-around-us:water}{rio}, uma declividade suave de \qty{10}{\%}, a mais
de \qty{1}{km} da \omterm{def:g7:earthquakes:fault}{falha}.

\textbf{11.} Por exemplo: travas diagonais e paredes de contraventamento
na estrutura, ligações fortes, ou isoladores embaixo do prédio; os
alunos treinam o que fazer num \omterm{def:g4:volcanoes-earthquakes-first-look:earthquake}{terremoto} (abaixar, proteger-se,
segurar-se).

\textbf{12.} \textbf{17 casas} de 28 ficam dentro da zona de \omterm{def:g5:water-shapes-the-land:flood}{inundação}:
$17 \div 28 \approx 0.61$, cerca de 3 casas em 5.
\end{solution}
